\chapter{Five sensors and the 165-band grid}\label{ch:sensors}

\begin{onething}
Five sensors hand you five incompatible things, and one decision fixes the worst of the
mismatch: resample every hyperspectral sensor onto one list of wavelengths, $10$ nm apart from
$460$ to $2{,}450$ nm with four atmospheric windows removed. That list has exactly $165$ entries,
and you can derive the $165$ on paper.
\end{onething}

\section{Why this chapter exists}\label{sec:sensors:motivation}

\Cref{ch:light} described a hyperspectral cube as if every sensor produced the same thing. They
do not. The five sensors this product accepts disagree about file format, about the order of the
axes in memory, about how stored integers become physical units, about which quality
information is shipped, and, most damagingly, about \emph{which wavelengths} they measure.

The last disagreement is the one that matters for a model. Band number $47$ of one sensor and band
number $47$ of another sit at different wavelengths. A network has one set of learned weights per
input position, so if you hand it both sensors unchanged, the weights for position $47$ have to
learn two contradictory things. The fix is a shared list of wavelengths, the
\textbf{common grid}, and the second half of this chapter derives it. The first half is the mess
that the fix hides, because every awkward corner of the ingestion code exists to hide it.

\section{The roster}\label{sec:sensors:roster}

\begin{table}[htbp]\centering
\caption{The five sensors. Band counts drift as products are reprocessed; the structural
differences do not. Sources: module docstrings of the dataset builders and
\texttt{docs/03-data-pipeline.md}.}\label{tab:sensors:roster}
\small
\begin{tabular}{@{}lllll@{}}\toprule
Sensor & Kind & Bands & Upload & Units after onboarding\\\midrule
PRISMA & hyperspectral & $239$ ($66 + 173$) & one \texttt{.he5} file & reflectance\\
EnMAP & hyperspectral & $224$ & a folder & reflectance\\
AVIRIS-NG & hyperspectral & about $425$ & a folder & reflectance\\
Landsat 9 & thermal & $1$ (band B10) & one \texttt{.tif} file & degrees Celsius\\
HotSat-1 & thermal & $1$ & a folder & raw DN\\\bottomrule
\end{tabular}
\end{table}

VNIR and SWIR are the names of a hyperspectral sensor's two detectors: one array for the visible
and near infrared, one for the shortwave infrared (\cref{tab:light:ranges}). PRISMA stores the
two as separate datasets inside one file, and its builder stitches them into one cube. For
AVIRIS-NG the fine-tuning shards use the $372$ bands its bad-band list marks good\tagM{docs/03-data-pipeline.md}.

\emph{File shape.} PRISMA and Landsat are single files; EnMAP, AVIRIS-NG and HotSat-1 are folders
in which the cube, the masks and the metadata are separate files. So ``the path to a scene''
sometimes means a file and sometimes a directory, resolved per sensor in
\file{backend/allotrope/sensors/source_path.py}.

\section{Three ways to lay a cube out in memory}\label{sec:sensors:layouts}

A computer's memory is one long line of numbers. A three-index array has to be flattened onto
that line, and the order in which the three indices are flattened is called its
\textbf{layout}\index{layout (array)}. Hyperspectral data uses three.

\begin{table}[htbp]\centering
\caption{The three cube layouts.}\label{tab:sensors:layouts}
\small
\begin{tabular}{@{}llp{0.22\linewidth}p{0.24\linewidth}@{}}\toprule
Layout & Axes & Contiguous in memory & Who uses it\\\midrule
BSQ, band-sequential & \shape{(C, H, W)} & one whole band image & EnMAP, Landsat, every model\\
BIL, band-interleaved-by-line & \shape{(H, C, W)} & one image row of one band & PRISMA (native)\\
BIP, band-interleaved-by-pixel & \shape{(H, W, C)} & one pixel's whole spectrum & visualisation, spectral matching\\\bottomrule
\end{tabular}
\end{table}

\subsection{Where an element lives: the offset formula}

NumPy, PyTorch and this codebase store arrays in \textbf{row-major}\index{row-major order}
order: the \emph{last} index changes fastest. For an array of shape \shape{(A, B, D)}, the
element at index $(a, b, d)$ sits at position
\begin{equation}\label{eq:sensors:offset}
\text{offset}(a, b, d) = a \cdot (B \cdot D) + b \cdot D + d
\end{equation}
counted from $0$. Read it as a mileage count: to reach element $(a,b,d)$ you first skip $a$ whole
blocks of size $B\cdot D$, then $b$ rows of size $D$, then $d$ single numbers. The three layouts
are three choices of which cube index plays $a$, $b$ and $d$:
\begin{align*}
\text{BSQ } (C, H, W):\quad &\text{offset} = c\,HW + i\,W + j,\\
\text{BIL } (H, C, W):\quad &\text{offset} = i\,CW + c\,W + j,\\
\text{BIP } (H, W, C):\quad &\text{offset} = i\,WC + j\,C + c.
\end{align*}

\begin{worked}[label={we:sensors:offset}]{One element, three addresses}
The cube of \cref{we:light:cube} has $C = 3$, $H = 2$, $W = 2$. Find where the element
``band $c=1$, row $i=0$, column $j=1$'' (the value $0.30$) lives in each layout. Do it on paper.

BSQ: $HW = 2 \times 2 = 4$, so
\[ \text{offset} = 1 \times 4 + 0 \times 2 + 1 = 4 + 0 + 1 = 5. \]
BIL: $CW = 3 \times 2 = 6$, so
\[ \text{offset} = 0 \times 6 + 1 \times 2 + 1 = 0 + 2 + 1 = 3. \]
BIP: $WC = 2 \times 3 = 6$, so
\[ \text{offset} = 0 \times 6 + 1 \times 3 + 1 = 0 + 3 + 1 = 4. \]
The same number lives at position $5$, $3$ or $4$ depending on the layout. Nothing in the file
says which layout it is; the reader has to know.
\end{worked}

\begin{worked}[label={we:sensors:misread}]{Reading BIL as if it were BSQ}
Write out the whole BIL memory line of that cube. BIL loops over rows $i$ (outer), then bands $c$,
then columns $j$ (inner). Row $0$: band $0$ gives $0.03, 0.05$; band $1$ gives $0.20, 0.30$;
band $2$ gives $0.14, 0.10$. Row $1$: $0.05, 0.03$; $0.30, 0.20$; $0.14, 0.10$. So memory holds
\[ 0.03,\ 0.05,\ 0.20,\ 0.30,\ 0.14,\ 0.10,\ 0.05,\ 0.03,\ 0.30,\ 0.20,\ 0.14,\ 0.10. \]
A reader who wrongly assumes BSQ takes the first $HW = 4$ numbers as band $0$:
\[ \text{``band 0''} = \begin{pmatrix} 0.03 & 0.05\\ 0.20 & 0.30\end{pmatrix}. \]
Its bottom row is really near-infrared data. Every ``band'' it builds is a mixture of rows of
different bands, and on a real scene the result looks like plausible but wrong static. If a cube
ever looks like structured noise, check the axis order before suspecting the data.
\end{worked}

The fix is a \textbf{transpose}\index{transpose (axes)}: reorder the axes without changing any
value. BIL \shape{(H, C, W)} to BSQ \shape{(C, H, W)} swaps the first two axes, in NumPy
\code{np.transpose(cube, (1, 0, 2))}. The codebase does it in
\code{ImageCubeOperations.convert_cube}
(\file{app/utils/image_transformation/image_cube_operations.py}) with a permutation of the axes.

The layouts are not arbitrary. BIP puts one pixel's whole spectrum next to itself in memory, so a
per-pixel spectral calculation is a fast linear walk; BSQ puts a whole band image together, so a
spatial filter over one band is. Choosing the wrong one costs nothing in correctness and a great
deal in speed.

\section{From stored integers to physical units}\label{sec:sensors:calibration}

Every sensor stores integers. There is no shared formula for turning them into physics.

\paragraph{EnMAP.} One uniform gain for all $224$ bands, set as \code{GAIN: float = 0.0001} in
the EnMAP reflectance transformer
\tagM{enmap_l2a_dn_to_surface_reflectance_transformer.py}:
\begin{equation}\label{eq:sensors:enmap}
\rho = \mathrm{DN} \times 0.0001.
\end{equation}

\paragraph{PRISMA.} A per-detector linear map read from the file's own metadata
(\file{app/utils/data_transformations/prs_l2d_dn_to_surface_reflectance_transformer.py}):
\begin{equation}\label{eq:sensors:prisma}
\rho = \frac{\mathrm{DN}}{65535}\,\bigl(\mathrm{L2ScaleMax} - \mathrm{L2ScaleMin}\bigr) + \mathrm{L2ScaleMin}.
\end{equation}
$65535 = 2^{16} - 1$ is the largest 16-bit unsigned integer, so $\mathrm{DN}/65535$ runs from $0$
to $1$, and the two scale attributes stretch that onto the reflectance range the provider chose
for this file.

\paragraph{Landsat 9.} The straight line of \cref{eq:light:landsat}, to Celsius.

\paragraph{HotSat-1.} No conversion: it ships raw 14-bit digital numbers around $5{,}000$ with a
spread of about $400$\tagM{app/models/training/inference_config.py}. \Cref{ch:bandpipeline} shows
what that does to a model fitted on other units.

\begin{worked}[label={we:sensors:dn}]{Two reflectance conversions}
\emph{EnMAP.} A pixel stores $\mathrm{DN} = 4{,}250$. By \cref{eq:sensors:enmap},
\[ \rho = 4{,}250 \times 0.0001 = 0.425. \]
That is $42.5\%$ reflectance, plausible for bright bare soil or a light roof; vegetation in red
light would be far lower, around $0.05$.

\emph{PRISMA.} Take $\mathrm{L2ScaleMin} = 0$ and $\mathrm{L2ScaleMax} = 0.7$ (illustrative
values\tagA; each real file carries its own). A pixel stores $\mathrm{DN} = 20{,}000$. Step by step:
\begin{align*}
\mathrm{DN}/65535 &= 20{,}000/65{,}535 \approx 0.30518,\\
\mathrm{L2ScaleMax} - \mathrm{L2ScaleMin} &= 0.7 - 0 = 0.7,\\
0.30518 \times 0.7 &\approx 0.21363,\\
\rho &= 0.21363 + 0 \approx 0.2136.
\end{align*}
\end{worked}

\section{Quality information, and reading bits}\label{sec:sensors:quality}

The sensors also disagree about which quality information they ship. EnMAP ships ready-made masks
for cloud, cirrus, haze, cloud shadow and snow, one \shape{(H, W)} layer each. PRISMA's L2D
product ships none. Landsat ships one integer layer, \code{QA_PIXEL}, in which each \emph{bit}
is a separate yes/no flag. That asymmetry is why the band-filter setting
\code{quality_masks_to_apply} exists and applies to EnMAP alone (\cref{ch:bandpipeline}).

\begin{definition}[Bit flag]\index{bit flag}
An integer written in base $2$ is a row of binary digits (bits), bit $k$ worth $2^k$. A
\textbf{bit flag} uses bit $k$ as a separate yes/no answer. To read bit $k$ of a value $q$, shift
$q$ right by $k$ places (divide by $2^k$ and drop the remainder) and keep the last bit (the
remainder on division by $2$). In code: \code{(q >> k) & 1}.
\end{definition}

The Landsat builder (\file{app/utils/dataset_builder/landsat_dataset_builder.py}, lines
187--207) marks a pixel unusable if any of bits $0$ (fill), $1$ (dilated cloud), $3$ (cloud),
$4$ (cloud shadow) or $15$ (high bit of the cirrus confidence) is set.

\begin{worked}[label={we:sensors:qa}]{Unpacking two QA values}
Take two example values\tagA, $q = 21{,}824$ and $q = 22{,}280$. First write $21{,}824$ as a sum of
powers of two, subtracting the largest power that fits each time:
\begin{align*}
21{,}824 - 16{,}384 &= 5{,}440 \quad (2^{14}),\\
5{,}440 - 4{,}096 &= 1{,}344 \quad (2^{12}),\\
1{,}344 - 1{,}024 &= 320 \quad (2^{10}),\\
320 - 256 &= 64 \quad (2^{8}),\\
64 - 64 &= 0 \quad (2^{6}).
\end{align*}
So bits $14, 12, 10, 8, 6$ are set. None of bits $0, 1, 3, 4, 15$ is set, so the pixel is
\emph{valid}. (Bit $6$ is the ``clear'' flag.)

Now $22{,}280$:
\begin{align*}
22{,}280 - 16{,}384 &= 5{,}896 \quad (2^{14}),\\
5{,}896 - 4{,}096 &= 1{,}800 \quad (2^{12}),\\
1{,}800 - 1{,}024 &= 776 \quad (2^{10}),\\
776 - 512 &= 264 \quad (2^{9}),\\
264 - 256 &= 8 \quad (2^{8}),\\
8 - 8 &= 0 \quad (2^{3}).
\end{align*}
Bit $3$ (cloud) is set, so the pixel is \emph{flagged}. Check bit $3$ the way the code does:
$22{,}280 \gg 3$ is $\lfloor 22{,}280/8 \rfloor = 2{,}785$, and $2{,}785$ is odd, so
$2{,}785 \mathbin{\&} 1 = 1$.
\end{worked}

\section{How the code hides the differences}\label{sec:sensors:helpers}

Two layers keep sensor-specific knowledge in one place.

\begin{itemize}
\item A \textbf{FileHelper}\index{FileHelper} wraps one physical file format and hands back raw
arrays and metadata: \code{HE5Helper} (PRISMA), \code{TIFHelper} (Landsat), \code{EnmapHelper},
\code{ENVIHelper} (AVIRIS-NG), \code{HotSatHelper}.
\item A \textbf{template}\index{template (file)} is a dictionary telling a helper where each
logical component lives inside the file, by one of three strategies
(\file{app/models/hyperspectral_concepts/references.py}): \code{FILE_REFERENCE} (an internal path
inside the container, such as PRISMA's SWIR cube dataset), \code{ROOT_METADATA_FIELD} (a root-level
attribute, such as PRISMA's list of SWIR centre wavelengths) and \code{DIRECT_PROPERTY_DEFINITION}
(a property on the opened dataset, such as rasterio's \code{crs}).
\end{itemize}
Templates are \emph{injected, not looked up}: the helper never contains a sensor-specific path,
which is what lets one contract serve HDF5, GeoTIFF and ENVI files. Only PRISMA, Landsat and
EnMAP have templates; AVIRIS-NG and HotSat-1 parse their own sidecar files, an inconsistency the
source notes honestly.

Three traps will each cost you an afternoon.
\begin{itemize}
\item \textbf{rasterio band indices start at 1.} Every other array in the codebase counts from
$0$. Asking rasterio for band $0$ is an error.
\item \textbf{PRISMA keeps its bad bands.} A band the sensor flagged unusable still occupies a
slice of the cube, with $\mathrm{FWHM} = 0.0$. Filter on the validity flag before trusting
anything about a band.
\item \textbf{The AVIRIS-NG sample cubes are not in the repository.}
\file{samples/aviris_samples/} holds the ENVI \code{.hdr} headers and \code{.sha256} checksums;
the \code{.bin} data files are absent, so the AVIRIS-NG path cannot run locally without them.
\end{itemize}

\section{Sensors that do not agree on wavelengths}\label{sec:sensors:disagree}

Now the disagreement that matters for a model. PRISMA, EnMAP and AVIRIS-NG all ``measure the
spectrum'', but at different wavelengths, with different spacing and different totals. In made-up
but representative numbers\tagA:
\begin{center}
\begin{tabular}{@{}ll@{}}
PRISMA & \dots, $1{,}000$, $1{,}012$, $1{,}024$, \dots\ nm\\
EnMAP & \dots, $1{,}003$, $1{,}011$, $1{,}019$, \dots\ nm\\
AVIRIS-NG & \dots, $1{,}001$, $1{,}006$, $1{,}011$, \dots\ nm
\end{tabular}
\end{center}
Training data is the bottleneck for a model. One model per sensor gives each model a third of the
data, triples the work of training and maintenance, and makes every new sensor a new model. One
model for all three gives it all the data, and lets a new sensor join without retraining, but only
if position $c$ of the input means the same wavelength for every sensor.

The fix is to agree, once, on a fixed list of wavelengths, and to \textbf{resample}%
\index{resampling} every sensor onto it before anything else touches the data. Resampling means:
``you measured at $1{,}001$ and $1{,}006$; I need a value at $1{,}000$; estimate one from your
neighbours''. The estimating is interpolation, derived in \cref{ch:bandpipeline}. This chapter
builds the list.

\begin{analogy}
Three weather stations log temperature at different times: one at 06:07, 06:19, 06:31; another at
06:03, 06:11, 06:19; a third every five minutes. To put all three on one chart you agree on the
times you will report (06:00, 06:10, 06:20, \dots) and estimate each station's temperature at
those times from its nearest readings. Two decisions remain: how far apart the reporting times
are (not finer than the slowest station's readings, or you invent detail), and which times to
leave out entirely (say, the hour each night when a station is switched off, because any value
reported there would be made up). The common grid makes exactly these two decisions for
wavelengths.
\end{analogy}

\section{Building the grid}\label{sec:sensors:grid}

The list is \code{DEFAULT_COMMON_WAVELENGTH_GRID} in \file{app/models/dataset/vendables.py},
built by this function (lines 35--62, docstring trimmed):
\begin{lstlisting}
def build_common_wavelength_grid(start_nm=460.0, end_nm=2450.0,
                                 spacing_nm=10.0, exclusion_ranges=None):
    if exclusion_ranges is None:
        exclusion_ranges = [
            (0, 450),       # low SNR
            (912, 978),     # water vapor
            (1131, 1152),   # water vapor
            (1350, 1450),   # water vapor absorption
            (1800, 1950),   # water vapor + CO2
        ]
    full_grid = np.arange(start_nm, end_nm + spacing_nm / 2, spacing_nm)  # (200,)
    mask = np.ones(len(full_grid), dtype=bool)
    for lo, hi in exclusion_ranges:
        mask &= ~((full_grid >= lo) & (full_grid <= hi))   # drop lo <= x <= hi
    return full_grid[mask]                                 # (165,)
\end{lstlisting}
Two decisions, then: a ladder of evenly spaced wavelengths, and a list of ranges cut out of it.
Each can be counted by hand.

\subsection{Step 1: the ladder, and the fencepost formula}

The ladder starts at $460$, steps by $10$, and stops at $2{,}450$:
\[ 460,\ 470,\ 480,\ \dots,\ 2{,}440,\ 2{,}450. \]
How many rungs? Number them: the rung at wavelength $x$ is rung number $k = (x - 460)/10$, so
$460$ is rung $0$, $470$ is rung $1$, and $2{,}450$ is rung $(2{,}450 - 460)/10 = 199$. Rungs
$0$ to $199$ are $200$ rungs. In general:
\begin{equation}\label{eq:sensors:fencepost}
\text{number of rungs} = \frac{\text{last} - \text{first}}{\text{step}} + 1.
\end{equation}
The ``$+1$'' is the \textbf{fencepost}\index{fencepost formula} correction: a fence $10$ metres long
with a post every metre has $11$ posts, because the division counts gaps, not posts.
\begin{align*}
\frac{2{,}450 - 460}{10} + 1 &= \frac{1{,}990}{10} + 1\\
&= 199 + 1\\
&= 200 \text{ rungs.}
\end{align*}

\begin{sourcenote}[: the half-step in \texttt{np.arange}]
\code{np.arange(start, stop, step)} stops \emph{before} \code{stop}, so
\code{np.arange(460, 2450, 10)} would end at $2{,}440$ and have $199$ entries. Adding half a
step makes the limit $2{,}455$, which includes $2{,}450$ without adding $2{,}460$. Verified:
\code{len(np.arange(460.0, 2450.0, 10.0))} is $199$ and the code's ladder has $200$ entries
(\file{research/textbook/tools/compute_sensors_output.txt}).
\end{sourcenote}

\subsection{Step 2: count the rungs inside each window}

The atmosphere is not transparent. At some wavelengths water vapour and carbon dioxide absorb the
light almost completely before it reaches the satellite; a band there measures noise, and noise
whose strength changes with humidity, so it looks like structure. Five ranges are removed, each
with both ends included (``$lo \le x \le hi$'', the last line of the loop).

Every rung is a multiple of $10$. The multiples of $10$ inside $[lo, hi]$ run from the first
multiple at or above $lo$ to the last at or below $hi$. Writing $\lceil y \rceil$ for ``round $y$
up'' and $\lfloor y \rfloor$ for ``round $y$ down'', those are $10\lceil lo/10\rceil$ and
$10\lfloor hi/10\rfloor$, and by \cref{eq:sensors:fencepost} with step $10$ (dividing through by
$10$):
\begin{equation}\label{eq:sensors:count}
\text{rungs removed} = \left\lfloor \frac{hi}{10} \right\rfloor - \left\lceil \frac{lo}{10} \right\rceil + 1,
\end{equation}
provided the range lies on the ladder (a range entirely below $460$ removes nothing). Do each
window on paper.

\begin{worked}[label={we:sensors:windows}]{Counting the five windows}
\emph{Window $(0, 450)$.} Every rung is at least $460 > 450$. Nothing removed: $0$.

\emph{Window $(912, 978)$.}
\begin{align*}
\lceil 912/10 \rceil &= \lceil 91.2 \rceil = 92,\\
\lfloor 978/10 \rfloor &= \lfloor 97.8 \rfloor = 97,\\
97 - 92 + 1 &= 6.
\end{align*}
The six rungs are $920, 930, 940, 950, 960, 970$. Check the edges: $910 < 912$ is kept;
$980 > 978$ is kept.

\emph{Window $(1131, 1152)$.}
\begin{align*}
\lceil 113.1 \rceil &= 114, \qquad \lfloor 115.2 \rfloor = 115,\\
115 - 114 + 1 &= 2 \quad (1{,}140 \text{ and } 1{,}150).
\end{align*}

\emph{Window $(1350, 1450)$.} Both ends land on rungs and both are included:
\begin{align*}
\lceil 135 \rceil &= 135, \qquad \lfloor 145 \rfloor = 145,\\
145 - 135 + 1 &= 11 \quad (1{,}350, 1{,}360, \dots, 1{,}450).
\end{align*}

\emph{Window $(1800, 1950)$.}
\begin{align*}
180 \text{ and } 195: \quad 195 - 180 + 1 &= 16 \quad (1{,}800, \dots, 1{,}950).
\end{align*}

\emph{Total.}
\begin{align*}
\text{removed} &= 0 + 6 + 2 + 11 + 16 = 35,\\
\text{remaining} &= 200 - 35 = 165.
\end{align*}
\end{worked}

\begin{sourcenote}[: 165 checked against the source]
Running the code's own mask over the ladder gives the same per-window counts $0, 6, 2, 11, 16$
and \code{len(DEFAULT_COMMON_WAVELENGTH_GRID)} $= 165$, first value $460.0$, last $2{,}450.0$
(\file{research/textbook/tools/compute_sensors.py} and its saved output)\tagC. Nothing about
$165$ is chosen directly: it is a consequence of the ladder and the five windows.
\end{sourcenote}

\begin{keypoint}
$165 = 200 - (0 + 6 + 2 + 11 + 16)$: a $10$ nm ladder from $460$ to $2{,}450$ nm minus the rungs
the atmosphere ruins.
\end{keypoint}

\subsection{The five segments}

The surviving rungs fall into five unbroken runs, $10$ nm apart within each run, with jumps
between runs. Count each with \cref{eq:sensors:fencepost}:
\begin{align*}
460 \text{ to } 910:&\quad (910 - 460)/10 + 1 = 45 + 1 = 46,\\
980 \text{ to } 1{,}130:&\quad (1{,}130 - 980)/10 + 1 = 15 + 1 = 16,\\
1{,}160 \text{ to } 1{,}340:&\quad (1{,}340 - 1{,}160)/10 + 1 = 18 + 1 = 19,\\
1{,}460 \text{ to } 1{,}790:&\quad (1{,}790 - 1{,}460)/10 + 1 = 33 + 1 = 34,\\
1{,}960 \text{ to } 2{,}450:&\quad (2{,}450 - 1{,}960)/10 + 1 = 49 + 1 = 50.
\end{align*}
Check: $46 + 16 + 19 + 34 + 50 = 165$. The four jumps (from $910$ to $980$, $1{,}130$ to $1{,}160$,
$1{,}340$ to $1{,}460$, $1{,}790$ to $1{,}960$) are the four windows that touch the ladder.
\Cref{fig:sensors:grid} draws the whole grid to scale.

\begin{figure}[htbp]\centering
\begin{tikzpicture}[x=0.0066cm,y=1.5cm]
  % windows, hatched, to scale
  \foreach \lo/\hi in {912/978,1131/1152,1350/1450,1800/1950}{
    \fill[pattern=north east lines,pattern color=maskgrey] (\lo,-0.05) rectangle (\hi,0.75);
    \draw[maskgrey] (\lo,-0.05) rectangle (\hi,0.75);}
  \fill[pattern=north east lines,pattern color=maskgrey] (400,-0.05) rectangle (450,0.75);
  \draw[maskgrey] (400,-0.05) -- (450,-0.05) -- (450,0.75) -- (400,0.75);
  % the 200 rungs
  \foreach \w in {460,470,...,2450}{
    \pgfmathtruncatemacro{\ex}{or(or(and(\w>=912,\w<=978),and(\w>=1131,\w<=1152)),or(and(\w>=1350,\w<=1450),and(\w>=1800,\w<=1950)))}
    \ifnum\ex=1 \draw[maskgrey,line width=0.5pt] (\w,0.05) -- (\w,0.45);
    \else \draw[spec,line width=0.6pt] (\w,0) -- (\w,0.6); \fi}
  % axis
  \draw[->] (400,-0.15) -- (2520,-0.15);
  \foreach \t in {460,1000,1500,2000,2450}{\draw (\t,-0.15) -- (\t,-0.25) node[below,font=\scriptsize]{\t};}
  \node[font=\scriptsize] at (1460,-0.75) {wavelength (nm)};
  % family braces
  \foreach \a/\b/\n in {460/910/46,980/1130/16,1160/1340/19,1460/1790/34,1960/2450/50}{
    \draw[decorate,decoration={brace,amplitude=3pt},spec] (\a,0.68) -- (\b,0.68)
      node[midway,above=2pt,font=\scriptsize,spec]{\n};}
  % window labels
  \foreach \lo/\hi/\n in {912/978/6,1131/1152/2,1350/1450/11,1800/1950/16}{
    \node[font=\scriptsize,maskgrey!70!black] at ({(\lo+\hi)/2},1.25) {$-\n$};}
  \node[font=\scriptsize,maskgrey!70!black,anchor=east] at (450,1.25) {$0{-}450$: $-0$};
\end{tikzpicture}
\caption{The common grid drawn to scale. Ticks are the $200$ ladder rungs from $460$ to
$2{,}450$ nm; grey rungs fall inside the hatched atmospheric windows and are removed
($6 + 2 + 11 + 16 = 35$), teal rungs are kept. Braces give the number of kept bands in each of the
five segments ($46 + 16 + 19 + 34 + 50 = 165$). The window $0$--$450$ nm lies below the ladder and
removes nothing.}\label{fig:sensors:grid}
\end{figure}

\coursefig[0.95\linewidth]{iw-03-grid}{The same grid as drawn in the Indradhanu course (part 3):
a ladder of $200$ wavelengths loses $35$ rungs to four atmospheric windows, leaving $165$ bands in
five segments of $46$, $16$, $19$, $34$ and $50$.}{fig:sensors:coursegrid}

\begin{sourcenote}[: ``spectral families'' in the vendable]
The Indradhanu course (part 3) calls the five segments ``spectral families'' and says the vendable
records which of the five each band is in, in \code{spectral_family_order}. The source disagrees:
\code{spectral_family_order} holds values of the enum \code{SpectralFamily}
(\file{app/models/hyperspectral_concepts/spectral_family.py}), whose members are
\code{VNIR}, \code{SWIR} and \code{PANCHROMATIC}. It records which \emph{detector} a band came
from, not which of the five segments it lies in. On the common grid each band is labelled with
the detector of its nearest native band (\code{_assign_families} in
\file{app/utils/data_transformations/spectral_resampler.py}). The five segments have no name in
the code; they exist only as the jumps in the grid. \Cref{ch:bandpipeline} reports a problem with
that nearest-band labelling.
\end{sourcenote}

\subsection{From band index to wavelength and back}

Because the segments are evenly spaced inside, the wavelength of grid band $c$ ($0 \le c \le 164$)
is a piecewise formula. Inside a segment that starts at index $c_0$ and wavelength $\lambda_0$,
each step of the index adds $10$ nm:
\begin{equation}\label{eq:sensors:index}
\lambda(c) = \lambda_0 + 10\,(c - c_0), \qquad
(c_0, \lambda_0) \in \{(0, 460),\ (46, 980),\ (62, 1{,}160),\ (81, 1{,}460),\ (115, 1{,}960)\}.
\end{equation}
The segment starts are the running totals of the segment sizes: $0$, then $0 + 46 = 46$, then
$46 + 16 = 62$, then $62 + 19 = 81$, then $81 + 34 = 115$.

\begin{worked}[label={we:sensors:index}]{Which wavelength is band 43, and which band is 1,000 nm?}
\emph{Band 43.} $43 < 46$, so it is in the first segment, $(c_0, \lambda_0) = (0, 460)$:
\[ \lambda(43) = 460 + 10 \times (43 - 0) = 460 + 430 = 890 \text{ nm}. \]
\emph{1,000 nm.} It lies in the second segment ($980$ to $1{,}130$). Invert
\cref{eq:sensors:index}:
\begin{align*}
c &= c_0 + (\lambda - \lambda_0)/10\\
  &= 46 + (1{,}000 - 980)/10\\
  &= 46 + 2 = 48.
\end{align*}
\emph{Band 100.} $81 \le 100 < 115$, so $(c_0, \lambda_0) = (81, 1{,}460)$ and
\[ \lambda(100) = 1{,}460 + 10 \times 19 = 1{,}650 \text{ nm}. \]
All three agree with the code's array (\file{research/textbook/tools/compute_sensors_output.txt})\tagC.
\end{worked}

\begin{sourcenote}[: a numeric slip in the course]
The Indradhanu course (part 3) illustrates the ``wire format'' warning below with the phrase
``band 43 means 1000 nm''. On the real grid band $43$ is $890$ nm and $1{,}000$ nm is band $48$
(\cref{we:sensors:index}). The warning itself is correct; the example pair is not a real pair.
\end{sourcenote}

\section{Two design questions}\label{sec:sensors:why}

\subsection{Why 10 nm and not 5}

The grid has to respect the \emph{coarsest} contributing sensor. PRISMA's bands are about $12$ nm
wide (the code comment above \code{DEFAULT_COMMON_WAVELENGTH_GRID} says ``PRISMA \textasciitilde
12nm''). A value every $5$ nm would claim detail PRISMA never measured: the interpolator would
happily produce numbers, but they would be invented, and a model would learn to reproduce invented
structure. $10$ nm is fine enough to be useful and coarse enough that every grid value is
supported by real measurements.

\subsection{Why cut the windows out of the grid, rather than interpolate and discard later}

Suppose $1{,}350$ to $1{,}450$ nm were left in the grid. After band filtering there is no
real sensor data between $1{,}340$ and $1{,}460$ (\cref{ch:bandpipeline} drops those bands first).
An interpolator does not know that; asked for values at $1{,}350, 1{,}360, \dots, 1{,}450$, it
draws a line or curve from the last real sample below to the first real sample above, and returns
eleven values that were never measured. Those values would go into training, the model would learn
to reproduce them, and at inference a genuine feature in that region would be compared against
fiction. Removing the wavelengths from the grid means the interpolator is never asked the
question: the problem is designed out rather than cleaned up afterwards.

\begin{sourcenote}[: which interpolator resamples onto the grid]
The Indradhanu course (part 3) says the resampling onto the grid uses PCHIP. The source says
otherwise: \file{app/utils/data_transformations/spectral_resampler.py} uses \emph{linear}
interpolation, and its docstring explains why (``the native bands are densely sampled (6-12nm
apart) and the target grid is 10nm \dots\ PCHIP offers no practical quality gain but costs
significantly more memory''). PCHIP is used one step earlier, to fill gaps in partly valid
spectra (\file{spectral_interpolator.py}). The argument above holds for either interpolator: a
straight line across a $120$ nm gap is as invented as a curve. \Cref{ch:bandpipeline} derives both.
\end{sourcenote}

\section{The grid is a wire format}\label{sec:sensors:wire}

Nowhere in a saved model file does it say ``input position $48$ is $1{,}000$ nm''. The link is
purely positional: input position $c$ means whatever the $(c+1)$-th entry of the grid is. Three
things downstream are lined up with this exact list in this exact order:
\begin{itemize}
\item the model's input width, \code{in_channels: int = 165} in
\code{HyperspectralSegFormerMAEConfig} (\file{app/models/training/training_config.py}, line 88);
\item the $165$ means and $165$ standard deviations in \file{app/constants/hyperspectral.json}
(\cref{ch:bandpipeline});
\item every training shard, in which PRISMA and EnMAP patches are mixed and, at the tensor level,
indistinguishable.
\end{itemize}
AVIRIS-NG was added later and needed no model change, because it comes out of the band pipeline as
the same $165$-band cube; the resolver says so in a comment
(\file{backend/allotrope/foundation_models/resolver.py}, near line 259).

So if someone edited the exclusion ranges tomorrow, the model would still load (still $165$
inputs), the normalisation statistics would still load (still $165$ numbers), and every band would
silently refer to the wrong wavelength. No error, no warning, wrong answers. Treat the grid like a
database schema: changing it invalidates everything already stored against it.

\begin{drift}
\emph{Dated 2026-10-29.} The $165$-band grid is what Indradhanu v0.2.0 consumes. The model shipped
as Indradhanu V2 reads each sensor's \emph{native} bands instead: its resolver entry sets
\code{band_grid="native"} (\file{backend/allotrope/foundation_models/resolver.py}), and the
\code{band_filter_apply} option \code{use_common_wavelength_grid} is documented as ``Required by
Indradhanu (v0.2.0); Indradhanu V2 needs it false (native bands)''
(\file{backend/allotrope/action_types/band_filter_apply.py}). \Cref{ch:ceiling,ch:crossstem}
explain the model that makes this possible. The grid remains the default band-filter output and
the input of the classical detectors.
\end{drift}

\begin{sourcenote}[: two different default exclusion lists]
Two default lists of windows exist. \code{BandFilterConfig.exclusion_ranges} and
\code{build_common_wavelength_grid} use $(1350, 1450)$ and $(1800, 1950)$. The older
\code{DEFAULT_EXCLUSION_RANGES} in \file{app/utils/data_transformations/spectral_band_filter.py}
uses wider ones, $(1328, 1492)$ and $(1784, 1967)$, labelled ``Standard PRISMA exclusion
ranges''. The wider list is used only when \code{SpectralBandFilter} is built with
\code{exclusion_ranges=None}; every builder in this book passes
\code{band_filter_config.exclusion_ranges} explicitly, so the product uses the narrower list.
Separately, the description of \code{BandFilterConfig.common_wavelength_grid} in
\file{app/models/dataset/vendables.py} still says ``(450--2450nm, 201 bands)''; the grid the code
builds has $165$ bands starting at $460$ nm. The description is stale; the code wins.
\end{sourcenote}

\begin{confusions}
\confusion{Are the 165 bands evenly spaced?}{Within each of the five segments, yes, $10$ nm apart.
Across the whole list, no: there are four jumps. The list is ascending but not uniform.}
\confusion{Does resampling lose information?}{A little, always. That is the price of
comparability, and the grid is chosen (at $10$ nm) to keep the loss small.}
\confusion{Could we feed each sensor's native bands and pad to a fixed length?}{Then position $c$
would mean different wavelengths in different rows of one batch, which is exactly what stops one
fixed-position model from serving three sensors. (A model that is told each band's wavelength can
avoid this; \cref{ch:ceiling} builds one.)}
\confusion{The window list starts with $(0, 450)$, yet nothing is removed by it. Why keep it?}{It
guards the band \emph{filter}, which runs on native sensor bands, some of which lie below
$450$ nm. On the grid, which starts at $460$, it has no effect.}
\confusion{BIL, BSQ, BIP: is one of them correct?}{No. Each is fastest for a different access
pattern. The models all expect BSQ \shape{(C, H, W)}.}
\confusion{Are the two thermal sensors interchangeable?}{No. Landsat 9 arrives in degrees Celsius;
HotSat-1 arrives as raw DN around $5{,}000$. A model normalised for one sees the other as wildly
out of range (\cref{ch:bandpipeline}).}
\end{confusions}

\begin{checkyourself}
\item\label{ex:sensors:1} How many rungs does a ladder from $400$ to $2{,}500$ nm in steps of
$10$ nm have? How many in steps of $5$ nm? Use \cref{eq:sensors:fencepost} and show each line.
\item\label{ex:sensors:2} Suppose the window $(1131, 1152)$ were widened to $(1125, 1160)$. Use
\cref{eq:sensors:count} to count the rungs it removes, and give the new length of the grid.
\item\label{ex:sensors:3} Using \cref{eq:sensors:index}, find the wavelength of grid band $70$ and
of band $164$, and the index of $2{,}210$ nm.
\item\label{ex:sensors:4} A cube has shape $C = 4$, $H = 3$, $W = 5$. Compute the memory offset of
the element at band $2$, row $1$, column $3$ in BSQ, BIL and BIP layouts.
\item\label{ex:sensors:5} A Landsat \code{QA_PIXEL} value is $q = 21{,}840$. Write it as a sum of
powers of two and decide whether the builder marks the pixel valid. An EnMAP pixel stores
$\mathrm{DN} = 1{,}830$: what is its reflectance?
\end{checkyourself}

\nextchapter{ch:product}
